\section{Business Objective}

The business objective is to increase customer retention and revenue while reducing ineffective marketing costs. In practice, this means identifying which customers should be prioritized for retention marketing, so that limited marketing budget is spent on the customers most at risk of churning, rather than being spread evenly across the entire customer base.

\section{Business KPIs}

The following KPIs are used to evaluate the business impact of the project:

\begin{itemize}
    \item \textbf{Retention / Repurchase Rate} -- the proportion of customers who remain active (make another purchase) within the target window.
    \item \textbf{Revenue} -- total revenue generated from retained customers.
    \item \textbf{Marketing Cost} -- the cost of running a retention campaign on a targeted customer segment.
    \item \textbf{Cost per Retained Customer} -- marketing cost divided by the number of customers who actually stay active.
    \item \textbf{Estimated Profit / ROI} -- the net financial benefit of the campaign after accounting for marketing cost.
\end{itemize}

\section{Targeting Strategies Compared}

Three targeting strategies are compared to quantify the value added by modeling:

\begin{enumerate}
    \item \textbf{Random / No-model} -- customers are targeted without any prioritization, serving as the baseline.
    \item \textbf{RFM rule-based} -- customers are prioritized using simple Recency-Frequency-Monetary heuristics.
    \item \textbf{ML targeting} -- customers are prioritized using the predicted churn probability from a trained model.
\end{enumerate}
